\printedpage{29}
% head: A. Grothendieck : Dérivateurs, Ch. VII, Édition des Universités de Montpellier II et Paris VII
\msp{page 45}is described uniquely by saying that it is \textit{constant} on each of the ``connected components'' (in the sense of the numbering of the vertices of $I_\tau$) of the set $E' = E'(f) \subseteq I_\tau = [0,n]$ of vertices which are sources or targets of transition arrows degenerate for $f$, i.e. in $\mathrm{Tr}(\tau)\backslash E$. And on the $\alpha$-th element of $E$ ($\alpha \in [1, r = \operatorname{card} E]$), of index $i_\alpha$ in $\mathrm{Tr}(\tau)$, $f$ induces the functor
\[
\boldsymbol{\Delta}_{\varepsilon_\alpha} \simeq \underbrace{J_{i_\alpha}}_{\substack{\text{transition}\\ \text{interval}\\ \text{of index } i_\alpha\\ \text{in } I_\tau}} \longrightarrow I_{\tau'}
\]
given by \textit{the unique} \underline{Ty}-\textit{arrow} in $I_\tau$ from $J_{i_\alpha}$ to the $\alpha$-th factor, of type $\Delta^{n_\alpha}_{\varepsilon_\alpha}$, of $\tau'$, given by $(*)$
\[
\Delta^{n_\alpha}_{\varepsilon_\alpha} \simeq \text{type of } \boldsymbol{\Delta}^{\varepsilon_\alpha}_{n_\alpha}\ \ \ldots
\]

\subsection*{2$^o$ Determination of $\mathrm{Ob}\,(\uTy/\tau')$.}

If $(\tau, f) \in \mathrm{Ob}\,(\uTy/\tau')$, thus
\[
f : I_\tau \longrightarrow I_{\tau'},
\]
we first of all associate with it the subset
\[
\mathrm{Cr}(\tau', f) = E'_f = E' \subseteq I_{\tau'}
\]
consisting of the vertices (called ``critical for $f$'', cf. page 41) of $I_{\tau'}$ which are source or target of an arrow $f(u)$, where $u \in \mathrm{Tr}(\tau, f)$, i.e. $f(u)$ is not degenerate (\footnote[13]{isolate a lemma $f$ increasing unnecessary.}). $E'$ contains $f(I_{\tau\,\text{réd}})$ (i.e. the images under $f$ of the critical vertices of $I_\tau$),
\msp{page 46}but can be larger. (E.g. if $\tau = \tau'$, $f = \mathrm{id}$, $E' = I_{\tau'}$, which is thus strictly larger than $I_{\tau'\,\text{réd}}$ if $\tau'$ is not reduced.) For every $i \in E'$, let us consider $f^{-1}(\{i\}) \subseteq I_\tau$. Since $f$ \texttt{[is]} increasing, it is an interval of the totally ordered set $I_\tau$, so it is a subinterval of the interval category $I_\tau$. As such, it has a signature $\tau_i \in \Ty$. We thus have a map
\[
\eta_f \text{ or } \eta : E' \longrightarrow \Ty.
\]
I say that $(\tau, f) \in \mathrm{Ob}\,(\uTy/\tau')$ is uniquely determined by this pair
\[
(E', \eta) \begin{cases} E' \in \mathfrak{P}(I_{\tau'}) = \mathfrak{P}([0,n']) \\[1ex] \eta : E' \longrightarrow \Ty\ . \end{cases}
\]
Moreover, the latter is subject to the sole restriction that
\[
E' \supseteq \mathrm{Cr}(\tau') = \text{set of critical vertices of } I_{\tau'}
\]
\printedpage{30}
% head: A. Grothendieck : Dérivateurs, Ch. VII, Édition des Universités de Montpellier II et Paris VII
(where what amounts to the same thing, that $E'$ contains $0_{\tau'}$, $1_{\tau'}$ and that two consecutive elements of $E'$ are joined by an arrow of $I_{\tau'}$): % [sic] où -> ou, « ou ce qui revient au même »
\[
\left\{
\begin{array}{lcl}
\mathrm{Ob}\,(\Ty/\tau') & \simeq & \text{set of pairs } (E',\eta), \text{ with} \\[1ex]
 & & E' \subseteq [0, \underbrace{\ n'\ }_{=\,\mathrm{long}\,\tau'}], \quad E' \supseteq \mathrm{Cr}(\tau'), \\[1ex]
 & & \eta : E' \longrightarrow \Ty\ .
\end{array}
\right.
\]
Let
\[
0 = t_0 < t_1 < \ \cdots\ < t_r = n'
\]
be the sequence of elements of $E'$, let for \ovalbox{$i \in [1,N]$}
\[
\left\|
\begin{array}{l}
l_i = t_i - t_{i-1} \\[1ex]
\varepsilon_i \quad \text{signature of } [t_{i-1}, t_i] \text{ in } I_{\tau'}\ ,
\end{array}
\right.
\]
\msp{page 47}and for \ovalbox{$i \in [0,N]$}\,, let
\[
\|\ \ \tau_i = \eta(t_i).
\]
Then the type $\tau$ is given by
\[
\tau = \tau_0 \Delta_{\varepsilon_1} \tau_1 \Delta_{\varepsilon_2} \cdots \Delta_{\varepsilon_N} \tau_N, \tag*{$(*)$}
\]
\texttt{[of]} length $N + \sum_{i=0}^N \underbrace{\mathrm{long}(\tau_i)}_{\lambda_i \geq 0}$.

\subsection*{3$^o$ Description of $\mathrm{Fl}(\uTy)$.}

If
\[
f : \tau \longrightarrow \tau' \qquad \texttt{[is]} \ \text{in } \uTy\ ,
\]
we associate with it as before an integer \underline{$N \geq 0$} ($N = \operatorname{card}\mathrm{Cr}(\tau', f){-}1$, $N = 0$ if and only if $\tau' = 1$, type of the one-point interval), maps
\[
\left\{
\begin{array}{lcl@{\qquad\qquad}ccl}
{[1,N]} & \longrightarrow & [+1,-1] \times \mathbf{N}^*, & i & \longmapsto & (\varepsilon_i, l_i) \\[1ex]
{[0,N]} & \longrightarrow & \Ty & i & \longmapsto & \tau_i
\end{array}
\right.
\]
(\footnote[14]{\textbf{NB} If $N = 0$, $[1,N] = \emptyset$, and there are no $\varepsilon_i$, $l_i$; the map $\underbrace{[0,N]}_{=\,\{0\}} \longrightarrow \Ty$ reduces to the datum of $\tau_0 \in \Ty$.}). The arrow $f$ in $\uTy$ is determined by these data. The type $\tau$ is given by the formula $(*)$ above, the type $\tau'$ by % [sic] determinée -> déterminée, accent manquant
\[
\tau' = \Delta^{l_1}_{\varepsilon_1} \Delta^{l_2}_{\varepsilon_2}\ \cdots\ \Delta^{l_N}_{\varepsilon_N}, \qquad \text{length } n' = \textstyle\sum l_i, \tag*{$(*')$}
\]
\printedpage{31}
% head: A. Grothendieck : Dérivateurs, Ch. VII, Édition des Universités de Montpellier II et Paris VII
and the increasing map from $I_\tau = [0, N + \sum \mathrm{long}(\tau_i)]$ to $I_{\tau'} = [0, \sum l_i]$ is given thus:
\[
\begin{array}{l@{\qquad}cl}
\text{constant on } [0, 0+\lambda_0] & \longmapsto & 0 \\[1ex]
\text{constant on } [\lambda_0+1, \lambda_0+1+\lambda_1] & \longmapsto & l_1 \\[1ex]
\text{constant on } [\lambda_0+\lambda_1+2, \lambda_0+\lambda_1+2+\lambda_2] & \longmapsto & l_2 \\[1ex]
\ldots & & \\[1ex]
\text{constant on } [\lambda_0+\lambda_1+\cdots+\lambda_i+(i+1), \lambda_0+\cdots+\lambda_i+\lambda_{i+1}+(i+1)] & \longmapsto & l_i.
\end{array}
\]

\printedpage{32}
% head: A. Grothendieck : Dérivateurs, Ch. VII, Édition des Universités de Montpellier II et Paris VII
\msp{page 48}

\section*{6\quad Composition in $\uTy$}

We have a composition law (everywhere defined) in $\mathrm{Ob}\,(\uTy) = \Ty$, it is that of the free monoid on two generators $\Delta_+$, $\Delta_-$. This composition law is interpreted, at the categorical level (and no longer merely combinatorial) as an amalgamated sum
\[
I_{\tau\tau'} = I_\tau \amalg_{(e, \underline{1}_\tau, \underline{0}_{\tau'})} I_{\tau'},
\]
where $e$ is the unit category, and where $\underline{1}_\tau$, $\underline{0}_{\tau'}$ now represent the morphisms $e \longrightarrow I_\tau$, $e \longrightarrow I_{\tau'}$ , having as image the target element and the source element respectively of the oriented interval in question. We thus have a cocartesian square in (Cat)
\[
\begin{tikzcd}[column sep=small, row sep=small]
 & e \arrow[dl, hook', "\underline{1}_\tau"'] \arrow[dr, hook, "\underline{0}_{\tau'}"] & \\
I_\tau \arrow[dr, hook, "i_{\tau,\tau'}"'] & \text{\scriptsize cocart.} & I_{\tau'} \arrow[dl, hook', "j_{\tau,\tau'}"] \\
 & I_{\tau\tau'}, &
\end{tikzcd}
\]
where all the arrows are monic and fully faithful functors. That this diagram (obviously commutative) is indeed cocartesian in (Cat) amounts, by definition, to saying that for every $X$ in Cat, the transposed diagram
\[
\begin{tikzcd}[column sep=small, row sep=small]
 & X & \\
\uCh_\tau(X) \arrow[ur, "\beta_\tau" near end, "\text{\scriptsize (target)}" near start] & & \uCh_{\tau'}(X) \arrow[ul, "\sigma_{\tau'}"' near end, "\text{\scriptsize (source)}"' near start] \\
 & \uCh_{\tau\tau'}(X), \arrow[ul] \arrow[ur] &
\end{tikzcd}
\]

\texttt{[is cartesian]}

\msp{page 49}which amounts to the following tautology that a path $C$ in $X$ of type $\tau\tau'$ ``is nothing other'' than a pair $(c,c')$ of a path $c$ of type $\tau$ and a path $c'$ of type $\tau'$ ``put end to end'' (in this order), i.e. such that
\[
\mathrm{but}(c) = \mathrm{source}(c')
\]
(and we then have
\[
\begin{array}{rcll}
\mathrm{source}(C) & = & \mathrm{source}(c) & \\[1ex]
\mathrm{but}(C) & = & \mathrm{but}(c') & \qquad\quad ).
\end{array}
\]

\printedpage{33}
% head: A. Grothendieck : Dérivateurs, Ch. VII, Édition des Universités de Montpellier II et Paris VII
Here arises the everlasting problem whether we are going to denote this path by $c{\cdot}c'$, which at first sight seems to impose itself, or rather $c' \circ c$, regarding the path $c$ from $x = \mathrm{source}(c)$ to $y = \mathrm{but}(c)$ as a sort of morphism from $x$ to $y$ (cf. below the formalism of categories of fractions), $c'$ as a morphism from $y = \mathrm{source}(c')$ to $z = \mathrm{but}(c')$, and interpreting the composite $C$ as the composite of these two morphisms $c$, $c'$ (in this order from the point of view of the operations to be carried out), which is traditionally denoted $c' \circ c$ or $c' \cdot c$, and not $c \circ c'$ or $c \cdot c'$, alas! It is this notation that we are going to follow below. % [sic] traditionellement -> traditionnellement, orthographe

But for the moment we are not going to concern ourselves with paths in arbitrary categories $X$, but return to the formalism of the $\tau$ and the $I_\tau$ themselves. The categorical interpretation of $\tau\tau'$ as an amalgamated sum assures us that this composition operation is functorial:

\msp{page 50}
\[
\tau_1 \xrightarrow{\ f\ } \tau_2, \quad \tau'_1 \xrightarrow{\ f'\ } \tau'_2 \qquad\quad \text{gives } \tau_1\tau'_1 \xrightarrow{f*f'} \tau_2\tau'_2
\]
(we shall not write $f \cdot f'$ for the arrow denoted $f * f'$, because of the possible confusion with the composition of arrows in $\uTy$). We have transitivity in an obvious way: if moreover $\tau_2 \xrightarrow{\ g\ } \tau_3$, $\tau'_2 \xrightarrow{\ g'\ } \tau'_3$, \texttt{[we have]}
\[
(gf) * (g'f') = (g * g') \circ (f * f').
\]
Thus we have a composition bifunctor
\[
\uTy \times \uTy \longrightarrow \uTy \qquad\quad (c,c') \longmapsto c \cdot c' \text{ or } c * c'.
\]
(In fact, we also have a bifunctor
\[
\uTyp \times \uTyp \longrightarrow \uTyp,
\]
but we prefer now to forget the category \underline{Typ}, which gives a distinctly less simple formalism.)

This composition is \textit{strictly} associative and \textit{strictly unital} (the unit being of course the unit element 1 of the free monoid Ty, corresponding to the one-point category $\{0\}$), and not just ``up to isomorphism''. And with good reason, since in $\uTy$ the isomorphisms are identities. \textit{Thus} $\uTy$ \textit{is a monoid in the category} (Cat) of small categories!

\msp{page 49$'$ (il y a un décalage dans la numérotation)}This implies that not only $\mathrm{Ob}\,\uTy = \Ty$, but also $\mathrm{Fl}(\uTy)$, and more generally all the sets of the form $\mathrm{Hom}_{\mathrm{Cat}}(X, \uTy)$ are monoids, and this functorially in $X$, for example the components $\mathrm{Hom}_{\mathrm{Cat}}(\boldsymbol{\Delta}_i, \uTy)$ of the semi-simplicial set ``nerve'' of $\uTy$,
\[
\mathrm{Nerf}_i(\uTy) = \mathrm{Hom}_{\mathrm{Cat}}(\boldsymbol{\Delta}_i, \uTy),
\]
are monoids, and the semi-simplicial maps between them are monoid homomorphisms: the nerve of $\uTy$ is a \textit{semi-simplicial monoid.}